% !TeX root = equidistribution.tex
\section{Equidistribution}



\Yang{This section is fully by AI, need human to read and rewrite}


\subsection[status=draft,tag=0000000,label=\labelNFM{subsec:arithmetic-Equidistribution}]{Arithmetic equidistribution}

Let $X/K$ be a projective variety of dimension $d$ over a number field.
Recall that $(x_m) \subset X(\bar K)$ is \emph{generic} if
\[ \#\{m:x_m \in Z\} < \infty \qquad\text{for every proper closed }Z \subsetneq X. \]
For $v \in M_K$ and $x \in X(\bar K)$, put
\[ \delta_{x,v} \coloneqq \frac{1}{[K(x):K]}\sum_{\substack{\sigma:K(x)\hookrightarrow\bbC_v\\\sigma|_K=v}}\delta_{\sigma(x)}. \]
Here $K(x)$ is the residue field, and the sum is over embeddings extending the fixed embedding at $v$.

\begin{theorem}[label=\labelNFM{thm:equidistribution-of-small-points}, status=draft, tag=0000000]
Let $\bar L$ be a nef adelic line bundle with $L$ big, and let $(x_m) \subset X(\bar K)$ be generic.
Then
\[ h_{\bar L}(x_m) \longrightarrow h_{\bar L}(X) \quad\Longrightarrow\quad \delta_{x_m,v} \longrightarrow \mu_{\bar L,v} \qquad(v \in M_K), \]
weakly on $X_v^\an$.
See \cite[Theorem~5.4.3]{Yuan-Zhang-2026-AdelicLineBundles}.
\end{theorem}
\begin{proof}
\Yang{}
\end{proof}



























\subsection[status=draft,tag=0000000,label=\labelNFM{subsec:algebraic-Equidistribution}]{Algebraic equidistirbution in dynamics}

Let $K$ be a number field, $k = \bar K$, and $f:X = \bbP^d_k \to X$ an endomorphism of degree $q \geq 2$.
Put
\[ f^*\calO(1) = \calO(q), \qquad D \coloneqq \deg f = q^d. \]
Indeed, a curve in a fiber would give $0 = \deg(f^*\calO(1)|_C) = q\deg(C) > 0$.
Thus $f$ is finite; smoothness and the Cohen--Macaulay flatness criterion give flatness.
The formula for $D$ follows by intersecting $d$ pulled-back hyperplanes.
For integral $W$, write $\eta_W$ for its generic point.
We say that $C$ is totally invariant if $f^{-1}(C) = C$, and that $W$ belongs to a totally invariant cycle if $(f^p)^{-1}(W) = W$ for some $p \geq 1$.
The references are \cite[\S\S 2--5]{Gignac-2014-EquidistributionPreimages} and \cite[\S 7]{Gignac-2012-NoetherianSpaces}.

\begin{lemma}[label=\labelNFM{lem:algebraic-multiplicity-cocycle}, status=draft, tag=0000000]
For $x \in X$, $y = f(x)$, put
\[ e_f(x) \coloneqq \operatorname{length}_{\calO_{X,x}}(\calO_{X,x}/\frakm_y\calO_{X,x}), \qquad m_f(x) \coloneqq e_f(x)[\rkk(x):\rkk(y)]. \]
Then $e_f:X \to \{1,\ldots,D\}$ is upper semicontinuous, and
\[ e_{f^{n+m}} = e_{f^n}(e_{f^m}\circ f^n), \qquad \sum_{f(x)=y}m_f(x) = D. \]
For closed $x$, $m_f(x) = e_f(x)$; for integral $F$ with $f^p(F) = F$,
\[ [k(F):(f^p)^*k(F)] = q^{p\dim F}. \]
\end{lemma}
\begin{proof}
For a flat local map $(A,\frakm) \to B$ and an $\frakm$-primary ideal $\fraka$, tensor a composition series of $A/\fraka$ with $B$:
\[ \operatorname{length}_B(B/\fraka B) = \operatorname{length}_A(A/\fraka)\operatorname{length}_B(B/\frakm B). \]
This gives the cocycle identity; residue-field degrees multiply as well.
Since $f_*\calO_X$ is locally free of rank $D$,
\[ D = \dim_{\rkk(y)}\bigl((f_*\calO_X)\otimes\rkk(y)\bigr) = \sum_{f(x)=y}m_f(x). \]
At closed points both residue fields are $k$.

For upper semicontinuity, let $I$ be the diagonal ideal on $X\times_f X$, and $\pi$ the first projection.
The coherent sheaf
\[ \scrF \coloneqq \pi_*(\calO_{X\times_f X}/I^D) \]
has fiber dimension $e_f(x)$.
Indeed, after extending the fiber over $f(x)$ to $\rkk(x)$, separability singles out the diagonal Artinian factor with residue field $\rkk(x)$ and length $e_f(x)$.
On this factor $I^D = 0$; on every other factor $I$ is the unit ideal.
Thus $\dim_{\rkk(x)}\scrF(x) = e_f(x)$.
Fiber dimensions are upper semicontinuous by the minors of a presentation matrix.
Finally, for $r = \dim F$, the projection formula gives
\[ [k(F):(f^p)^*k(F)]\deg(F) = \bigl((f^p)^*\calO(1)\bigr)^r\cdot[F] = q^{pr}\deg(F). \]
\end{proof}

\begin{lemma}[label=\labelNFM{lem:noetherian-diagonal-limits-and-orbits}, status=draft, tag=0000000]
Every sequence of probability measures admits a subsequence converging on constructible sets to a positive finitely additive $\ell$, with $\ell(X) = 1$.
For $x \in X$ and bounded upper semicontinuous $\tau$,
\[ \frac1n\sum_{j<n}\tau(f^j(x)) \longrightarrow \frac1p\sum_{i<p}\tau(\eta_i), \]
where $\eta_i$ are the generic points of the components of the eventual orbit closure of $x$.
\end{lemma}
\begin{proof}
There are countably many closed sets, since their ideals are finitely generated over the countable field $k$.
Diagonal extraction gives convergence on closed sets; inclusion--exclusion gives convergence on constructible sets.
Positivity, finite additivity, and total mass one pass to the limit.

For forward averages, first suppose the eventual orbit closure $L$ is irreducible; discard finitely many initial terms.
Write
\[ d_N(A) \coloneqq N^{-1}\#(A\cap[0,N)), \qquad \bar d(A) \coloneqq \limsup_N d_N(A). \]
If a proper closed subset of $L$ is visited with positive upper density, choose a minimal such $W$.
Then $W$ is irreducible, and its proper closed subsets are visited with density zero.
Put $A = \{j:f^j(x) \in W\}$, $s = \bar d(A) > 0$, and $A-i = \{n \geq 0:n+i \in A\}$.
If $W\cap f^{-h}(W) \subsetneq W$ for every $h > 0$, then
\[ \bar d\bigl((A-i)\cap(A-j)\bigr) = 0 \qquad(i<j). \]
For each $M$, inclusion--exclusion gives
\[ M d_N(A) \leq 1+\sum_{0\leq i<j<M}d_N\bigl((A-i)\cap(A-j)\bigr)+O_M(N^{-1}). \]
Taking $d_{N_t}(A) \to s$ gives $Ms \leq 1$ for every $M$, a contradiction.
Hence, for some $h$,
\[ W \subset f^{-h}(W), \qquad f^h(W) = W. \]
The second equality follows from finiteness and preservation of dimension.
Once the orbit meets $W$, its tail lies in $\bigcup_{j<h}f^j(W) \subsetneq L$, contradicting density.
Thus every proper closed subset of $L$ is visited with density zero.

In general, the tail closures stabilize and $f(L) = L$.
Its components $L_0,\ldots,L_{p-1}$ form one cycle, since a single orbit is dense in $L$.
The subsequences modulo $p$ are dense in their respective components: the union of their closures is $L$.
Apply the preceding argument to $f^p|_{L_i}$.
The set $E_{i,\epsilon} = \{z \in L_i:\tau(z) \geq \tau(\eta_i)+\epsilon\}$ is proper and closed, so its visit density is zero; also $\tau\geq\tau(\eta_i)$ on $L_i$.
Hence, along that subsequence,
\[ 0 \leq \frac1N\sum_{j<N}\tau((f^p)^j(z))-\tau(\eta_i) \leq \epsilon+2\|\tau\|_\infty d_N(\{j:(f^p)^j(z)\in E_{i,\epsilon}\}). \]
Let $N \to \infty$, then $\epsilon \to 0$, and average the $p$ residue classes.
\end{proof}

\begin{lemma}[label=\labelNFM{lem:reverse-multiplicity-growth}, status=draft, tag=0000000]
The following limit exists and is upper semicontinuous:
\[ e_-(x) \coloneqq \lim_{n\to\infty}\left(\max_{f^n(y)=x}e_{f^n}(y)\right)^{1/n} = \max_{\substack{F\ni x\text{ integral, periodic}\\p=\operatorname{per}(F)}}e_{f^p}(\eta_F)^{1/p}. \]
\end{lemma}
\begin{proof}
Put
\[ \tau = \log e_f, \qquad S_n = \sum_{j<n}\tau\circ f^j, \qquad T_n(x) = \max_{f^n(y)=x}S_n(y). \]
For $c>0$, Noetherianity gives
\[ A_c \coloneqq \bigcap_{n\geq1}\{S_n\geq cn\} = \bigcap_{n=1}^M\{S_n\geq cn\}. \]
The functions $S_j/j$, $j\leq M$, take finitely many values.
Choose $0\leq a<c$ bounding their values below $c$.
An orbit segment avoiding $A_c$ splits into blocks of length at most $M$, each of average at most $a$, and a remainder of length less than $M$.
Thus, with $C=\|\tau\|_\infty$,
\[ S_n \leq an+MC. \]

Let $E_i$ be the components of $A_c$, and $L_i$ the eventual orbit closures of $\eta_{E_i}$.
For some common $s$, $f^s(E_i)\subset L_i$ for all $i$.
If $x\notin\bigcup_iL_i$, a reverse orbit ending at $x$ avoids $A_c$ before its last $s$ terms; hence
\[ T_n(x) \leq an+MC+sC, \qquad \limsup_n T_n(x)/n < c. \]
If $x\in L_i$, the preceding lemma gives a generic-point cycle average $\beta_i\geq c$, since $\eta_{E_i}\in A_c$.
Lift $x$ backwards inside this cycle.
A lift in a component $F$ satisfies $\tau(y)\geq\tau(\eta_F)$, so
\[ \liminf_n T_n(x)/n \geq \beta_i \geq c. \]

Put $c=\limsup_nT_n(x)/n$.
If $c=0$, the limit is zero and is attained at $\eta_X$, since $e_f(\eta_X)=1$.
If $c>0$, the two implications put $x$ in some $L_i$ and force the liminf to equal $c$.
Every periodic generic point specializing to $x$ gives the same backward-lifting lower bound; one of them attains $c$.
This proves the maximum formula.
For each $c>0$,
\[ \left\{x:\lim_nT_n(x)/n\geq c\right\} = \bigcup_iL_i. \]
These sets are closed; for $c\leq0$ the level set is $X$.
Exponentiate to obtain $e_-$.
\end{proof}

\begin{lemma}[label=\labelNFM{lem:ramification-gap-and-fiber-count}, status=draft, tag=0000000]
For integral $W\subset X$, put $r=\dim W$, $c=d-r$.
Then
\[ e_-(\eta_W)\leq q^c, \qquad e_-(\eta_W)=q^c \iff (f^p)^{-1}(W)=W\text{ for some }p\geq1, \]
and, for closed $z$,
\[ \#(f^{-n}(z)\cap W)\leq\deg(W)q^{nr}. \]
The union $\calE_f$ of the proper totally invariant closed subsets is itself proper, closed, and totally invariant.
\end{lemma}
\begin{proof}
For periodic integral $F\supset W$, $f^p(F)=F$,
\[ e_{f^p}(\eta_F)=\frac{m_{f^p}(\eta_F)}{q^{p\dim F}}\leq q^{p\operatorname{codim}F}. \]
Equality holds exactly when $\eta_F$ is its only preimage.
Flat going-down, after base change to $F$, makes every component of $(f^p)^{-1}(F)$ dominate $F$.
Thus equality is equivalent to $(f^p)^{-1}(F)=F$.
The maximum formula gives $e_-(\eta_W)\leq q^c$; equality forces $\operatorname{codim}F=c$, hence $F=W$.

If $z\notin f^n(W)$ the fiber is empty.
Otherwise take $r$ general hyperplanes through $z$, meeting $f^n(W)$ in finitely many points.
Their pullbacks contain $f^{-n}(z)\cap W$, and B\'ezout gives
\[ \#(f^{-n}(z)\cap W)\leq\bigl((f^n)^*\calO(1)\bigr)^r\cdot[W]=q^{nr}\deg(W). \]
For $r=0$ this is immediate.

Put $B_c=\{e_-\geq q^c\}$.
If $W$ has codimension $c$ and belongs to a totally invariant cycle, a component $F$ of $B_c$ containing $W$ satisfies
\[ q^c\leq e_-(\eta_F)\leq q^{\operatorname{codim}F}\leq q^c. \]
Thus $F=W$.
Each $B_c$, $1\leq c\leq d$, is proper and has finitely many components, since $e_-(\eta_X)=1$.
There are only finitely many such $W$; their union is proper and totally invariant.
Conversely, the components of any totally invariant closed set form such cycles: $f$ permutes them and flatness excludes extra generic preimages.
This union is therefore $\calE_f$.
\end{proof}

\begin{lemma}[label=\labelNFM{lem:gignac-escape-from-proper-subset}, status=draft, tag=0000000]
Put $\nu_{n,a}=D^{-n}(f^n)^*\delta_a$, with $a$ closed.
If $\nu_{n_i,a}\to\ell$ on constructible sets and $W$ is a minimal closed subset with $\ell(W)>0$, then $(f^p)^{-1}(W)=W$ for some $p\geq1$.
\end{lemma}
\begin{proof}
Minimality makes $W$ irreducible.
Suppose $W$ does not belong to a totally invariant cycle; put $c=\operatorname{codim}W$, $s=\ell(W)>0$.
The preceding lemma and upper semicontinuity give a nonempty open $U\subset W$ with
\[ e_-|_U<q^c, \qquad \ell(U)=\ell(W)=s. \]
We claim that some $z\in f^{-n_I}(a)\cap U$ satisfies
\[ \limsup_i\nu_{n_i-n_I,z}(W)>0. \]
If not, set $n_0=0$ and
\[ R(z,I)\coloneqq\limsup_i\nu_{n_i-n_I,z}(U), \qquad z\in f^{-n_I}(a). \]
Then $R(z,I)=0$ for $z\in U$.
Whenever $R(z,I)\geq s$, choose $J>I$ with $\nu_{n_J-n_I,z}(U)>s/2$.
The multiplicity cocycle gives
\[ \nu_{n_i-n_I,z}=\sum_{f^{n_J-n_I}(w)=z}\nu_{n_J-n_I,z}(\{w\})\nu_{n_i-n_J,w}. \]
Taking mass on $U$ and the limsup, the terms with $w\in U$ vanish:
\[ R(z,I)\leq(1-s/2)\max_{w\notin U}R(w,J). \]
Starting at $R(a,0)=s$, iteration gives values at least $s(1-s/2)^{-j}>1$ for large $j$, a contradiction.
This proves the claim.

Choose $e_-(z)<\lambda<q^c$.
For all large $h$, every $y\in f^{-h}(z)$ has $m_{f^h}(y)=e_{f^h}(y)\leq\lambda^h$.
Therefore, with $r=d-c$,
\[ \nu_{h,z}(W)\leq D^{-h}\lambda^h\deg(W)q^{rh}=\deg(W)(\lambda/q^c)^h\longrightarrow0, \]
contradicting the claim.
\end{proof}

\begin{theorem}[label=\labelNFM{thm:gignac-algebraic-preimages}, status=draft, tag=0000000, title={Gignac}]
Let $f:\bbP^d_{\bar K}\to\bbP^d_{\bar K}$ have degree $q\geq2$.
For $a\in\bbP^d(\bar K)\setminus\calE_f$ and proper closed $Z\subsetneq\bbP^d_{\bar K}$,
\[ \nu_{n,a}(Z)=q^{-dn}\sum_{y\in f^{-n}(a)\cap Z}m_{f^n}(y)\longrightarrow0. \]
See \cite[Theorem~5.1]{Gignac-2014-EquidistributionPreimages}.
\end{theorem}
\begin{proof}
Otherwise take $\nu_{n_i,a}(Z)\geq\epsilon>0$ and a diagonal limit $\ell$ on constructible sets.
Choose a minimal closed $W\subset Z$ with $\ell(W)>0$.
The preceding lemma gives a totally invariant cycle whose union $C$ contains $W$.
Then
\[ \ell(C)>0\quad\Longrightarrow\quad f^{-n}(a)\cap C\neq\varnothing\text{ for some }n\quad\Longrightarrow\quad a\in C\subset\calE_f, \]
a contradiction.
\end{proof}

% \subsection[status=draft,tag=0000000,label=\labelNFM{subsec:distribution-of-periodic-points}]{Distribution of periodic points}

Here $k$ may be any algebraically closed field.
Let $f:\bbP^d_k\to\bbP^d_k$ have degree $q\geq2$.
Put
\[ S_n=\operatorname{Fix}(f^n)(k), \qquad N_n=1+q^n+\cdots+q^{dn}, \]
\[ R_n(Z)=\{y\in S_n:f^j(y)\notin Z\text{ for }0\leq j<n\}. \]
We will use the condition
\[ \operatorname{Fix}(f^n)\text{ is reduced for every }n\geq1. \tag{R}\label{eq:reduced-fixed-point-hypothesis} \]
Equivalently, $\det(Df^n_y-I)\neq0$ for all $y\in S_n$.

\begin{theorem}[label=\labelNFM{thm:algebraic-periodic-avoidance}, status=draft, tag=0000000]
Assume (R), and let $\varnothing\neq Z\subsetneq\bbP^d$ have dimension $r$.
For some $C_Z>0$ independent of $n$,
\[ 0\leq N_n-\#R_n(Z)\leq nC_Z(q^n+1)^r, \qquad \#R_n(Z)=q^{dn}(1+o(1)), \]
\[ \mu_{S_n}(Z)\coloneqq\frac{\#(S_n\cap Z)}{N_n}\leq\frac{C_Z(q^n+1)^r}{N_n}\longrightarrow0. \]
\end{theorem}
\begin{proof}
By \hyperref[lem:algebraic-fixed-point-counts]{the counting lemma in the Green-current section}, reducedness gives $\#S_n=N_n$ and $\#(S_n\cap Z)\leq C_Z(q^n+1)^r$; for reducible $Z$, sum over its components.
The map $f$ permutes $S_n$, with inverse $f^{n-1}$, so
\[ \#(S_n\setminus R_n(Z))\leq\sum_{j<n}\#\{y\in S_n:f^j(y)\in Z\}=n\#(S_n\cap Z). \]
Now $r<d$ gives
\[ \frac{N_n}{q^{dn}}\longrightarrow1, \qquad \frac{n(q^n+1)^r}{q^{dn}}\longrightarrow0. \]
Dividing the counting bounds proves the assertions.
For $Z=\varnothing$, they are immediate.
\end{proof}

\begin{remark}[label=\labelNFM{rmk:reducedness-is-an-extra-hypothesis}, status=draft, tag=0000000]
Condition (R) is extra: $f(z)=z+z^2$ has a double fixed point at $0$.
The power map $[x_i]\mapsto[x_i^q]$ satisfies (R) in characteristic zero, since in a fixed-point chart
\[ x_i^{q^n}-x_i=0, \qquad (x_i^{q^n}-x_i)'\in\{-1,q^n-1\}. \]
Without (R), B\'ezout bounds distinct points, not their ambient fixed-point multiplicities.
\end{remark}

\begin{theorem}[label=\labelNFM{thm:dinh-zhong-periodic-avoidance-cited}, status=draft, tag=0000000, title={Dinh--Zhong, cited}]
For $f:\bbP^d_\bbC\to\bbP^d_\bbC$ of degree $q\geq2$ and every proper closed $Z$,
\[ \#R_n(Z)=q^{dn}(1+o(1)). \]
Here points are distinct and their periods divide $n$.
\end{theorem}
This is \cite[Theorem~1.1]{Dinh-Zhong-2023-PeriodicPointsFujiki}: the dynamical degrees are $1,q,\ldots,q^d$.
The general theorem remains a cited input; under (R), the preceding algebraic proof suffices.




































\subsection[status=draft,tag=0000000,label=\labelNFM{subsec:dynamical-equidistribution}]{Complex-dynamical equidistribution}

Let $f:\bbP^d_K\to\bbP^d_K$ have degree $q\geq2$, with $K\subset\bbC$ a number field.
Write $\mu_f=\mu_{\bar L_f,v}=c_1(\bar L_f)_v^d$ for the canonical measure at the chosen complex embedding.
For a finite Galois-invariant cycle $P_n=\sum_y m_n(y)[y]$, put
\[ M_n=\sum_y m_n(y), \qquad \mu_{P_n}=M_n^{-1}\sum_y m_n(y)\delta_y. \]

\begin{theorem}[label=\labelNFM{thm:arithmetic-equidistribution-of-dynamical-cycles}, status=draft, tag=0000000]
For $a\in\bbP^d(\bar K)\setminus\calE_f$,
\begin{gather*}
q^{-dn}\sum_{y\in f^{-n}(a)}m_{f^n}(y)\delta_y\longrightarrow\mu_f,\\
q^{-dn}\sum_{f^n(y)=y}i_n(y)\delta_y\longrightarrow\mu_f,
\end{gather*}
where $i_n(y)$ is the graph--diagonal intersection multiplicity at $(y,y)$.
\end{theorem}
\begin{proof}
Enlarge $K$ so that $a$ is $K$-rational.
For backward cycles, Gignac's theorem and the height identity give
\[ M_n=q^{dn}, \qquad \mu_{P_n}(Z)\longrightarrow0, \qquad \widehat h_f(y)=q^{-n}\widehat h_f(a)\quad(y\in P_n). \]
For periodic cycles, $M_n=N_n$ and $\widehat h_f(y)=0$.
Use algebraic avoidance under (R), and the cited Dinh--Zhong theorem otherwise:
\[ 0\leq\sum_{y\in Z\cap S_n}i_n(y)\leq N_n-\#R_n(Z)=o(q^{dn}). \]
Thus both cases satisfy
\[ \max_{y\in\operatorname{Supp}P_n}\widehat h_f(y)\longrightarrow0, \qquad \mu_{P_n}(Z)\longrightarrow0\quad(Z\subsetneq\bbP^d\text{ closed}). \]
The \hyperref[claim:equidistirbution-from-arithmetic-equidistirbution]{Galois-orbit criterion in the Green-current section} gives $\mu_{P_n}\to\mu_f$.
Finally $M_n/q^{dn}\to1$.
\end{proof}
